% ==================== 1.3 现存难点及问题分析 ====================
\section{现存难点及问题分析}
\label{sec:challenges}

通过对国内外研究现状的分析，本节梳理月面巡视器自主行走面临的核心难点，凝练需要解决的关键科学问题。

\subsection{月面巡视器自主行走的核心难点}
\label{subsec:core_difficulties}

月面巡视器自主行走面临环境、本体、感知、决策等多维度的技术挑战，具体表现在以下方面：

\textbf{（1）月面环境复杂多变}

月面环境具有极端性和不确定性双重特征。极端性表现在：昼夜温差高达300\degree C以上，重力仅为地球的1/6，月尘具有强附着性和磨蚀性。不确定性表现在：月壤力学参数空间变化大，地形起伏不可预测，光照条件随太阳高度角持续变化。这些因素导致巡视器的行走性能难以准确预测，增加了自主决策的难度。

\textbf{（2）轮-壤相互作用机理不清}

月壤具有独特的物理力学特性：颗粒棱角分明、内聚力低、密度随深度变化。轮-壤相互作用涉及复杂的接触力学、散体力学和摩擦学问题，现有理论模型难以准确描述巡视器车轮与月壤的动态交互过程，导致动力学预测存在较大误差。

\textbf{（3）感知信息不完整}

受传感器性能、光照条件、月尘干扰等因素影响，巡视器获取的环境信息往往不完整、不准确。暗影区域光照不足导致视觉感知困难，月尘遮挡降低相机成像质量，通信中断造成数据传输延迟。这些因素制约了巡视器的环境理解和态势感知能力。

\textbf{（4）自主决策能力不足}

现有巡视器的自主决策能力有限，主要依赖地面规划与控制。在通讯受限条件下，巡视器难以独立应对复杂地形、动态障碍等突发情况。自主决策面临的挑战包括：如何快速生成可行路径、如何平衡效率与安全、如何处理规划与执行的偏差等。

\textbf{（5）虚实交互协同困难}

数字孪生技术要求数字空间与物理空间保持实时同步和双向交互。然而，地-月通讯时延、传感器噪声、模型误差等因素导致虚拟实体与物理实体之间存在偏差。如何实现虚实之间的精确映射和协同优化，是数字孪生技术应用的核心难题。

\subsection{三大科学问题提出}
\label{subsec:scientific_problems}

基于上述核心难点的分析，本文提出面向月面巡视器自主行走数字孪生技术研究的三大科学问题：

\textbf{科学问题一：月面巡视器自主行走数字孪生系统如何构建？}

数字孪生系统的构建是应用数字孪生技术的基础。月面巡视器自主行走数字孪生系统涉及物理实体（巡视器）、虚拟实体（数字模型）、服务功能（感知、规划、决策）、孪生数据（环境、状态、行为数据）和连接（数据传输、指令交互）五个维度。如何针对月面环境的特殊性和巡视器自主行走的需求，对经典五维模型进行适应性增强？如何构建物理实体与虚拟实体之间的映射关系？如何实现数据的实时采集、传输和处理？这些问题构成了月面巡视器自主行走数字孪生系统构建的核心科学问题。

\textbf{科学问题二：如何实现月面环境与巡视器动力学的精确数字孪生？}

精确的环境和动力学模型是数字孪生发挥作用的前提。月面环境包括地形、月壤、光照、温度等多物理场要素，这些要素具有强耦合、强时变的特点。巡视器动力学涉及轮-壤相互作用、多体运动、滑移特性等复杂力学过程。如何建立高保真、高效率的月面环境数字孪生模型？如何构建考虑轮-壤相互作用的巡视器动力学数字孪生模型？如何通过模型校准和参数辨识提高模型精度？这些问题构成了环境和动力学数字孪生建模的核心科学问题。

\textbf{科学问题三：如何在数字孪生环境下实现巡视器的智能感知与决策？}

智能感知与决策是数字孪生的核心价值所在。在数字孪生环境下，巡视器可以利用虚拟空间中的计算资源进行复杂的环境分析、行为预测和策略优化。如何在数字孪生平台上实现月面岩石等障碍目标的精确识别？如何设计基于数字孪生的路径规划与避障策略？如何实现虚拟决策向物理执行的可靠转化？这些问题构成了数字孪生环境下智能感知与决策的核心科学问题。

上述三大科学问题相互关联、层层递进，构成了本文研究的逻辑主线。科学问题一解决"系统架构"问题，科学问题二解决"模型基础"问题，科学问题三解决"功能实现"问题。
